\chapter{Image Processing}
\label{ch:images}

A digital image is a grid of numbers. Each number records how much light fell on
one small square of a sensor, and everything an image-processing system does is
arithmetic on that grid. Blurring averages neighbours. Edge detection takes
differences. Thresholding compares every pixel with one number, and morphology
takes the largest or smallest value in a window. The mathematics is old and well
understood: convolution, the Fourier transform, order statistics, and the
geometry of sets. What makes it hard in practice is scale. A modest photograph
has a quarter of a million pixels, and a pipeline may touch each of them dozens
of times.

This chapter shows how Mathilda handles images. It starts with the
representation, because a surprising number of image bugs come from misreading
it. It then works through the standard toolbox: geometry, tone and colour,
convolution and filtering, derivatives and edges, thresholding, and mathematical
morphology. It ends with volumetric images and a set of short worked tasks in the
style of a cookbook. Each operation is shown twice. A small matrix, printed at
the prompt, lets you check the arithmetic pixel by pixel. A larger image, rendered
as a figure, shows what the operation looks like.

Three themes run through the chapter. The first is that an image is an
ordinary expression, so the tools of the rest of the book work on it and on its
pixels. The second is checkable mathematics. Most operations obey an exact
identity (duality, idempotence, conservation of the mean, a reflected kernel),
and the sessions test those identities directly. The third is speed. Mathilda
keeps pixels in machine buffers, and the operations were benchmarked against
SciPy and scikit-image. The measurements are quoted where they matter and
collected in Section~\ref{sec:img-speed}.

Mathilda does not ship a library of sample photographs, so every image in this
chapter is built from a formula. That has one advantage: you can type any input
line yourself and get exactly the picture printed here.

\section{An image is an expression}
\label{sec:img-expr}

\B{Image} wraps a rectangular array of numbers. The simplest way to make a test
image is to write down a formula for the brightness at each pixel and let
\B{Table} evaluate it. A \emph{zone plate}\index{zone plate} is the classic
choice: rings whose spatial frequency rises steadily with distance from the
centre, so a single picture contains every frequency from zero up to the sampling
limit.

\plotcell[0.5\linewidth]{image-processing/zone}

The image is \(160\) pixels wide and \(128\) tall. It is deliberately not square,
for a reason the next session makes clear.

\subsection{Pixels, types, and the transposition trap}
\label{ssec:img-types}

A byte image stores integers from \(0\) to \(255\). The accessors report its
shape, its type, and its pixels:

\begin{codepairs}
\pair{image-processing/basics/1}{Two rows of three pixels, written as bytes.}
\pair{image-processing/basics/2}{\B{ImageQ} confirms that the data made a valid
image.}
\pair{image-processing/basics/3}{\B{ImageDimensions} gives the width first,
while \B{ImageData} is an array of rows. The two orders are transposed.}
\pair{image-processing/basics/4}{\B{ImageType} was inferred from the values: all
integers, all within \(0..255\).}
\pair{image-processing/basics/5}{With its type as the second argument,
\mcode{ImageData} returns the stored integers unchanged.}
\pair{image-processing/basics/6}{Without it, every pixel is scaled into
\([0,1]\). The byte \(255\) becomes exactly \(1.0\).}
\pair{image-processing/basics/7}{\B{FullForm} shows the canonical form
\mcode{Image[data, type]}. The pixels sit in an \B{NDArray}, a machine
buffer.}
\end{codepairs}

\index{image!dimensions}The transposition in line 3 is the most common source
of silently wrong image code. The part
\mcode{ImageData[img][[y, x]]} is row \(y\), counted from the top, and column
\(x\), counted from the left, so the array is height by width.
\mcode{ImageDimensions} follows the convention of a screen and reports width
first. Mathematica uses the same pair of conventions, and so does Mathilda.

\calloutnote{Pitfall}{A square test image cannot tell the two orders apart. Code
that confuses \mcode{\{width, height\}} with \mcode{Dimensions[ImageData[img]]}
passes every test on a \(64\times64\) image and fails on the first photograph.
Build your test images non-square, as this chapter does.}

\index{image!type}The type fixes the range of a stored value. Mathilda infers it
from the values alone, which makes it predictable:

\begin{center}
\small
\begin{tabular}{@{}ll@{}}
\toprule
data & inferred type \\
\midrule
all integers, every value in \(\{0, 1\}\) & \mcode{"Bit"} \\
all integers, every value in \(0..255\) & \mcode{"Byte"} \\
anything else & \mcode{"Real"} \\
\bottomrule
\end{tabular}
\end{center}

\begin{codepairs}
\pair{image-processing/types/1}{The same shape of data written three ways gives
three types. A real \(0.\) and an integer \(0\) are different data.}
\pair{image-processing/types/2}{A bit image scales to \(0.0\) and \(1.0\).}
\pair{image-processing/types/3}{Ragged rows are not an image, so \mcode{Image}
stays unevaluated and \mcode{ImageQ} says so.}
\pair{image-processing/types/4}{Neither is a string.}
\pair{image-processing/types/5}{Real data is kept as given, even outside
\([0,1]\). Clamping happens only when the image is written to a file.}
\end{codepairs}

\calloutnote{Under the hood}{\texttt{src/image.c} validates the data once, when
\mcode{Image} first normalises to its two-argument canonical form. A builtin that
cannot evaluate returns its input unchanged, so without a canonical form a valid
\mcode{Image[data]} and a malformed \mcode{Image["hello"]} would look alike.
Because the canonical form is only produced by the validator, the accessors
recheck only that the rows have equal length, which costs \(O(\text{height})\).
The first implementation revalidated every pixel on every call, and
\mcode{ImageDimensions} of a \(512\times512\) image took \(0.59\) ms. It now takes
\(0.58\,\mu\)s (\texttt{docs/spec/builtins/image-processing.md}).}

\subsection{Colour channels}
\label{ssec:img-colour}

A colour image adds a third axis. Each pixel is a list of channel values, so the
array is height by width by channels, interleaved:

\begin{codepairs}
\pair{image-processing/colour/1}{A \(2\times2\) image of red, green, blue and
white.}
\pair{image-processing/colour/2}{\B{ImageChannels} is \(3\). The data has a
third axis of length \(3\).}
\pair{image-processing/colour/3}{Row 1, column 2 is the green pixel.}
\pair{image-processing/colour/4}{\B{ColorConvert} to grey applies the Rec.~601
luminance weights.}
\pair{image-processing/colour/5}{The weights sum to one, so white stays
white.}
\pair{image-processing/colour/6}{The result has one channel.}
\end{codepairs}

\index{luminance!Rec.~601}The grey value of a colour pixel is
\[
  Y = 0.299\,R + 0.587\,G + 0.114\,B .
\]
Green dominates because the eye is most sensitive to it. Every Mathilda
operation that needs one brightness from a colour pixel (thresholding, edge
detection, template matching) uses these same weights, so the operations agree
with each other about what ``brightness'' means.

A channel can be pulled out with \B{Part} and viewed as an image in its own
right. Here red rises from left to right and green from top to bottom:

\plotcell[0.8\linewidth]{image-processing/channels}

\section{Making, saving, and loading images}
\label{sec:img-io}

\B{Export} writes an image to PNG, JPEG, or BMP according to the file extension,
and \B{Import} reads PNG, JPEG, BMP, GIF, TGA, PSD, HDR, and PNM files back.
\B{RandomImage} makes noise, which is the standard input for testing a filter:

\begin{codepairs}
\pair{image-processing/io/1}{\mcode{Export} returns the file name, so it can feed
\mcode{Import} directly.}
\pair{image-processing/io/2}{Read the file back.}
\pair{image-processing/io/3}{The samples survive exactly. The type is
\mcode{"Real"}, because \mcode{Import} scales 8-bit samples by \(1/255\).}
\pair{image-processing/io/4}{Uniform noise on \([0,1]\), four wide and three
high.}
\pair{image-processing/io/5}{The pixels.}
\pair{image-processing/io/6}{\B{SeedRandom} makes the noise reproducible. It
uses the same stream as \B{RandomReal}.}
\pair{image-processing/io/7}{With no arguments, \mcode{RandomImage} gives a
\(150\times150\) grey image.}
\end{codepairs}

The option \mcode{ColorSpace -> "RGB"} gives three independent channels:

\plotcell[0.5\linewidth]{image-processing/random}

\calloutnote{Under the hood}{\texttt{src/imageio.c} decodes and encodes with the
public-domain \texttt{stb\_image} headers, vendored so that \mcode{Import}
cannot fail for lack of a system library. They are compiled with internal
linkage in that one file. That keeps them from colliding with the copy of
\texttt{stb\_image} inside Raylib, which caused a link failure before the change.
JPEG output uses a fixed quality of \(90\). PNG is lossless and is the right
choice for results you intend to measure.}

\section{Geometry}
\label{sec:img-geometry}

Geometric operations move pixels without changing their values.
\B{ImageReflect}, \B{ImageRotate}, \B{ImagePad}, and \B{ImageCrop} are easiest to
understand on a small non-square image whose pixels are all different:

\begin{codepairs}
\pair{image-processing/geometry/1}{Pixels \(0.1\) to \(0.6\), in two rows.}
\pair{image-processing/geometry/2}{\mcode{ImageReflect} with no direction flips
top to bottom.}
\pair{image-processing/geometry/3}{\mcode{Left} (or \mcode{Right}) flips left to
right.}
\pair{image-processing/geometry/4}{A half turn reverses both axes.}
\pair{image-processing/geometry/5}{A quarter turn swaps width and height.}
\pair{image-processing/geometry/6}{Four quarter turns give back the identical
expression. \B{SameQ} tests identity, with no tolerance.}
\pair{image-processing/geometry/7}{\mcode{ImagePad} with a number pads with
zeros.}
\pair{image-processing/geometry/8}{\mcode{"Fixed"} replicates the edge pixel
instead.}
\pair{image-processing/geometry/9}{Cropping back to the original size undoes the
padding exactly.}
\pair{image-processing/geometry/10}{With no size, \mcode{ImageCrop} trims a
uniform border. The border colour is read from a corner.}
\end{codepairs}

The same operations on a picture, an asymmetric ramp, show which way each one
moves the content. From the left: the original, the vertical flip, the
horizontal flip, and the half turn.

\plotcell[\linewidth]{image-processing/reflect}

Padding needs a rule for the values it invents. \mcode{ImagePad} offers three: a
constant (zero by default), \mcode{"Fixed"}, which copies the nearest edge pixel
outward, and \mcode{"Reflected"}, which mirrors the image across its edge:

\plotcell[0.9\linewidth]{image-processing/pad}

\calloutnote{Theory}{The padding rule is the boundary condition of every
neighbourhood operation. A filter at the edge of an image reads pixels that do
not exist, and it has to be told what they are. Zero padding darkens the border
of a blurred image, which looks like lens vignetting. Replicating the edge keeps a
constant image constant under any smoothing filter, and Mathilda's filters all use
that rule. \mcode{ImagePad[img, m, "Fixed"]} makes the same rule available
explicitly, so padding first and filtering second agrees with filtering alone.}

\subsection{Rotation by an arbitrary angle}

A quarter turn sends every pixel to another pixel's position. A rotation by any
other angle does not, because a rotated grid does not land on the original grid.
\mcode{ImageRotate} then \emph{resamples}\index{resampling!rotation}. For each
output pixel it rotates the position backwards, finds where that point falls in
the source, and interpolates bilinearly between the four surrounding pixels.
Points that fall outside the source are black:

\plotcell[0.8\linewidth]{image-processing/rotate}

The third panel rotates the square by \(\pi/4\) and back again. The square
returns, slightly softened by two rounds of interpolation, but the corners of the
frame do not. They left the image on the first rotation and came back as black.

\calloutnote{Under the hood}{\texttt{src/imagegeom.c} keeps right angles on a
separate path that only permutes indices. Nothing is interpolated there, so four
quarter turns are the identity under \mcode{===}. Sending a quarter turn through
the bilinear resampler would give the same pixels only if the half-pixel
convention were exactly right. The general path uses \emph{inverse mapping}: it
loops over destination pixels and asks where each came from. Forward mapping, which
loops over the source and writes where each pixel lands, leaves holes wherever
the rotation stretches the grid.}

\subsection{Resizing and aliasing}
\label{ssec:img-resize}

\B{ImageResize} changes the number of pixels. Enlarging is interpolation.
Shrinking is harder, because it throws information away, and the way it is thrown
away matters. A fine checkerboard shows the problem:

\begin{codepairs}
\pair{image-processing/resize/1}{A \(4\times4\) checkerboard of \(0\) and
\(1\).}
\pair{image-processing/resize/2}{Halving each side by area averaging gives a flat
\(0.5\), the average brightness.}
\pair{image-processing/resize/3}{Point sampling keeps every other pixel. Those are
all \(0\), so the pattern turns black.}
\pair{image-processing/resize/4}{A small image to enlarge.}
\pair{image-processing/resize/5}{\mcode{"Nearest"} enlargement repeats each
pixel.}
\pair{image-processing/resize/6}{The default for enlarging is bilinear
interpolation.}
\pair{image-processing/resize/7}{A single number sets the width. The height
follows the aspect ratio.}
\pair{image-processing/resize/8}{A \(60\times40\) noise image.}
\pair{image-processing/resize/9}{Reducing it to \(40\times25\), a factor of
\(1.5\) one way and \(1.6\) the other, keeps the mean brightness exactly.}
\end{codepairs}

\begin{theorem}[Nyquist--Shannon]
\index{sampling theorem}\index{aliasing}
A signal sampled at spacing \(h\) is determined by its samples only if it
contains no frequency above \(1/(2h)\). A component above that limit is not
lost. It reappears in the samples as a lower frequency, its \emph{alias}.
\end{theorem}

Shrinking an image by a factor \(s\) multiplies the sample spacing by \(s\), so
every frequency above the new limit must be removed \emph{before} the new samples
are taken. Point sampling removes nothing, and the zone plate shows the result.
Below, the original is on the left. In the middle it has been reduced \(4\times\)
by point sampling and enlarged again for display. On the right the reduction used
area averaging:

\plotcell[\linewidth]{image-processing/aliasing}

The fine outer rings of the middle panel have turned into false patterns: whole
new sets of circles centred in the corners. They are aliases. Area averaging
replaces each new pixel by the mean of the source region it covers, which is a box
low-pass filter and a resample in one step. The high frequencies average out to
grey and no false structure appears.

\calloutnote{Under the hood}{\mcode{ImageResize} uses area averaging whenever
either axis shrinks and bilinear interpolation when both grow. The area weights
use fractional coverage, so a \(3\to2\) reduction is as correct as \(4\to2\), and
every source pixel contributes the same total weight. That is why the mean is
preserved exactly in line~9. Coordinates are centre-aligned: output pixel \(i\)
samples source position \((i+\tfrac12)s - \tfrac12\). The simpler map \(i s\)
shifts the image by half a pixel at every scale except \(1{:}1\).}

\calloutnote{Performance}{On a \(512\times512\) image, area reduction to
\(256\times256\) takes \(0.63\) ms and bilinear enlargement to \(1024\times1024\)
takes \(2.76\) ms, against \(2.0\) ms and \(14.3\) ms for scikit-image
(\texttt{docs/spec/builtins/image-processing.md}, ``Storage: machine buffers'';
benchmark sources in \texttt{benchmarks/64-image-resize}).}

\section{Tone: point operations and the histogram}
\label{sec:img-tone}

A \emph{point operation}\index{point operation} replaces each pixel by a function
of its own value. It ignores the neighbours, so it cannot blur or sharpen. It can
change brightness, contrast, and the distribution of grey levels.
\B{ImageAdjust} with no parameters stretches the range so the darkest pixel
becomes \(0\) and the brightest \(1\). With a list \(\{c, b, g\}\) it applies
contrast \(c\), brightness \(b\), and gamma \(g\) by the curve
\[
  v' = \Bigl(\operatorname{clip}_{[0,1]}\bigl((v - \tfrac12)(1 + c) + \tfrac12 + b\bigr)\Bigr)^{1/g}.
\]

\begin{codepairs}
\pair{image-processing/tone/1}{A dim, flat image: all values between \(0.40\)
and \(0.55\).}
\pair{image-processing/tone/2}{The stretch maps the range onto \([0,1]\)
linearly.}
\pair{image-processing/tone/3}{Stretching again changes nothing. The stretch is
idempotent.}
\pair{image-processing/tone/4}{Contrast \(1\) doubles distances from mid-grey.
Mid-grey itself stays fixed.}
\pair{image-processing/tone/5}{Brightness adds a constant.}
\pair{image-processing/tone/6}{Gamma \(2\) takes square roots and lifts the
shadows.}
\end{codepairs}

\calloutnote{Pitfall}{The contrast, brightness, and gamma curve is Mathilda's own
documented formula. Mathematica does not publish the exact curve its
\mcode{ImageAdjust} uses, so parametric adjustments may differ from Mathematica's
in the last digits or more. The plain stretch has no such ambiguity.}

The gamma parameter bends the ramp without moving its ends. From the top: gamma
\(\tfrac12\), the untouched ramp, gamma \(2\), and contrast \(+1\), which
saturates both ends:

\plotcell[0.7\linewidth]{image-processing/gamma}

\subsection{The histogram, and equalising it}

\B{ImageLevels} returns the histogram as data, a list of \(\{\text{level},
\text{count}\}\) pairs. The result is data you can compute with:

\begin{codepairs}
\pair{image-processing/tone/7}{A bit image has two natural levels. Here two
pixels are \(0\) and four are \(1\).}
\pair{image-processing/tone/8}{With \(4\) bins, each pixel counts at the nearest
of the levels \(0, \tfrac13, \tfrac23, 1\).}
\pair{image-processing/tone/9}{\B{HistogramTransform} spreads four distinct
values evenly over \([0,1]\).}
\end{codepairs}

\index{histogram equalisation}The last line is \emph{histogram equalisation}. If
\(F\) is the cumulative distribution of the pixel values, so \(F(v)\) is the
fraction of pixels no brighter than \(v\), the transform replaces each value \(v\)
by \(F(v)\). The result is as close to uniformly distributed as the data allows.
Four equally common values go to \(\tfrac14, \tfrac12, \tfrac34, 1\) regardless of
where they started.

The difference between a stretch and equalisation shows on an image whose values
are crowded into a narrow band. Below, from the top: the original, the stretch,
and the equalised image. The stretch is linear, so the disk and the background
keep their relative contrast. Equalisation gives each grey level screen space in
proportion to how many pixels use it. The large background gets most of the
range, and the fine banding shows the \(256\) bins.

\plotcell[0.4\linewidth]{image-processing/tone}

\calloutnote{Under the hood}{\texttt{src/imagecolor.c} computes the equalising
map from the Rec.~601 luminance and applies it to a colour image as a ratio,
multiplying every channel of a pixel by the same factor. Equalising each channel
separately would change the balance between them and shift the hue. The
\mcode{ImageAdjust} stretch takes \(0.76\) ms on a \(512\times512\) image against
\(7.7\) ms for scikit-image's \mcode{rescale\_intensity}, and \mcode{ImageLevels}
takes \(0.66\) ms against \(1.6\) ms for NumPy's \mcode{histogram}.}

\section{Colour}
\label{sec:img-colourops}

\mcode{ColorConvert} reduces a colour image to grey. The three primaries show the
Rec.~601 weights directly: pure green is the brightest, pure blue the darkest.

\plotcell[0.7\linewidth]{image-processing/grayscale}

\subsection{Quantisation by median cut}

\B{ColorQuantize} reduces an image to at most \(n\) colours. It uses
\emph{median cut}\index{median cut}\index{colour quantisation}. Think of every
pixel as a point in the RGB cube. Start with one box around all of them. Repeatedly
take the box whose widest single channel spans the largest range, and split it at
the median of that channel. When there are \(n\) boxes, replace every pixel by the
mean colour of its box.

\begin{codepairs}
\pair{image-processing/colourops/1}{Two clusters of grey, split into two colours.
Each pair collapses to its mean.}
\end{codepairs}

On the colour ramp below, two colours split the red axis, four split red and
green, and sixteen begin to follow the diagonal of the blue channel:

\plotcell[\linewidth]{image-processing/quantize}

\calloutnote{Theory}{\(k\)-means clustering gives a palette with lower total
error than median cut, but it starts from random centres. Its palette would depend
on the state of the random generator, so it could not be tested or documented as
a fixed result. Median cut is deterministic. Splitting the box with the widest
spread, rather than the box with the most pixels, puts palette entries where the
colours actually differ.}

\subsection{Replacing colours, and the alpha channel}

\B{ColorReplace} replaces every pixel within a Euclidean distance of a colour.
The default tolerance is \(0.02\):

\begin{codepairs}
\pair{image-processing/colourops/2}{Pure red, an almost-red, blue, and white.}
\pair{image-processing/colourops/3}{Only pure red is replaced. The almost-red
pixel is \(0.042\) away, outside the default tolerance.}
\pair{image-processing/colourops/4}{A tolerance of \(0.05\) catches it too.}
\end{codepairs}

\plotcell[0.5\linewidth]{image-processing/colorreplace}

With several rules, the nearest target colour wins, so the answer does not depend
on the order in which the rules are listed.

An \emph{alpha channel}\index{alpha channel} records opacity. \B{AlphaChannel}
reads it, \B{SetAlphaChannel} attaches one, and \B{RemoveAlphaChannel} either
drops it or composites the image over a background:

\begin{codepairs}
\pair{image-processing/colourops/5}{An image without alpha is fully opaque, so
its alpha channel is all ones.}
\pair{image-processing/colourops/6}{A white image at half opacity.}
\pair{image-processing/colourops/7}{Grey plus alpha is two channels.}
\pair{image-processing/colourops/8}{Composited over black, half-transparent white
is grey.}
\pair{image-processing/colourops/9}{Dropping the alpha forgets the transparency
and leaves white.}
\end{codepairs}

\calloutnote{Pitfall}{\mcode{ColorConvert} currently converts only to
\mcode{"Grayscale"} (or \mcode{"Gray"}). Other colour spaces such as
\mcode{"HSB"} and \mcode{"LAB"} leave the call unevaluated.}

\section{Convolution}
\label{sec:img-convolution}

Most filters are linear. Each output pixel is a weighted sum of the input pixels
around it, with the same weights everywhere. Such an operation is a
\emph{convolution}\index{convolution}. For an image \(f\) and a kernel \(k\) whose
centre is at index \((0, 0)\),
\[
  (f * k)[y, x] = \sum_{i, j} k[i, j]\, f[y - i,\, x - j].
\]
The minus signs matter. They mean the kernel is \emph{reflected} before it is slid
across the image. The same sum with plus signs is a
\emph{correlation}\index{correlation}. \B{ImageConvolve} computes the first and
\B{ImageCorrelate} the second:

\begin{codepairs}
\pair{image-processing/convolve/1}{A single bright pixel, a \emph{delta}.}
\pair{image-processing/convolve/2}{Convolving a delta reproduces the kernel. The
delta response of a filter \emph{is} its kernel.}
\pair{image-processing/convolve/3}{Correlation gives the kernel reversed.}
\pair{image-processing/convolve/4}{A constant image convolved with a kernel that
sums to \(1\) stays constant, border included.}
\pair{image-processing/convolve/5}{A kernel summing to \(3\) triples the
brightness.}
\pair{image-processing/convolve/6}{The result is always a real image, even for a
byte input.}
\end{codepairs}

Convolution and correlation agree on any symmetric kernel, such as a Gaussian or a
box, which is why confusing them is easy. They differ on an asymmetric kernel, as
lines 2 and 3 show. Line 4 depends on the border rule. Every filter in Mathilda
replicates the edge pixel, the \mcode{"Fixed"} rule of \mcode{ImagePad}, so a
constant image has constant neighbours everywhere.

\calloutnote{Theory}{Convolution is the natural operation because it commutes
with translation: shifting the image and then filtering gives the same result as
filtering and then shifting. By the convolution theorem, convolution becomes
pointwise multiplication after a Fourier transform,
\(\widehat{f * k} = \hat f \cdot \hat k\). A blur is therefore a filter that
multiplies high spatial frequencies by small numbers, and \(\hat k\) is its
frequency response.}

\subsection{Gaussian and box kernels}

The two standard smoothing kernels are the Gaussian and the box.
\B{GaussianMatrix}\mcode{[r]} gives a \((2r+1)\times(2r+1)\) Gaussian with
standard deviation \(r/2\), normalised to sum to one. \B{BoxMatrix}\mcode{[r]} is a
matrix of ones, unnormalised as in Mathematica. \B{GaussianFilter} and
\B{MeanFilter} are the corresponding filters:

\begin{codepairs}
\pair{image-processing/gauss/1}{A radius-1 Gaussian. The centre carries most of
the weight.}
\pair{image-processing/gauss/2}{It sums to exactly one at any radius.}
\pair{image-processing/gauss/3}{The box is all ones, so it brightens by a factor
of nine unless you divide.}
\pair{image-processing/gauss/4}{A delta in a \(5\times5\) image.}
\pair{image-processing/gauss/5}{Filtering the delta reproduces
\mcode{GaussianMatrix[1]}, up to rounding in the last bit.}
\pair{image-processing/gauss/6}{\mcode{GaussianFilter} is the same computation
as \mcode{ImageConvolve} with \mcode{GaussianMatrix}, bit for bit.}
\pair{image-processing/gauss/7}{\mcode{MeanFilter} spreads the delta evenly over
the \(3\times3\) window.}
\pair{image-processing/gauss/8}{It agrees exactly with the normalised box
kernel.}
\end{codepairs}

\calloutnote{Pitfall}{Normalising a truncated Gaussian by the analytic constant
\(1/(2\pi\sigma^2)\) leaves its sum slightly below one, because the tails beyond
the kernel's edge are missing. Every pass through such a filter would darken the
image a little. \mcode{GaussianMatrix} divides by the sum of the entries it
actually has, which is why line 2 prints exactly \(1.0\).}

Increasing the radius removes finer and finer detail. On the zone plate the outer
rings, which carry the highest frequencies, go first. From the left: radius \(0\)
(no change), \(8\), \(16\), and \(32\).

\plotcell[\linewidth]{image-processing/gaussian}

\subsection{Separability}
\label{ssec:img-separable}

A direct \((2r+1)\times(2r+1)\) convolution costs \((2r+1)^2\) multiplications per
pixel. For most smoothing kernels that is far more than necessary.

\begin{definition}
\index{separable kernel}
A kernel \(k\) is \emph{separable} if it is an outer product,
\(k[i, j] = u[i]\, v[j]\), which is to say a matrix of rank one.
\end{definition}

A separable convolution factors into a pass down the columns with \(u\) followed
by a pass along the rows with \(v\). The cost falls from \((2r+1)^2\) to
\(2(2r+1)\) per pixel: at radius \(8\), from \(289\) multiplications to \(34\).
The Gaussian is separable because
\(e^{-(x^2+y^2)/2\sigma^2} = e^{-x^2/2\sigma^2} e^{-y^2/2\sigma^2}\). So is the box.

\calloutnote{Under the hood}{\mcode{ImageConvolve} tests \emph{every} kernel for
rank one, whatever its origin. It pivots on the entry of largest magnitude, so a
kernel with zero corners (every derivative kernel has them) still factors, and it
checks each entry against the product to a relative tolerance of \(10^{-12}\). A
Gaussian factors to about \(10^{-16}\), so real separable kernels pass easily,
while a kernel of rank two fails by many orders of magnitude. Treating a rank-two
kernel as separable would give a completely different filter. Putting the test in
\mcode{ImageConvolve}, which \mcode{GaussianFilter} calls, is what keeps line 6
above an identity under \mcode{===}.}

\calloutnote{Performance}{At \(512\times512\), a separable Gaussian takes
\(1.48\) ms at radius \(2\) and \(3.44\) ms at radius \(8\). SciPy's general
\mcode{ndimage.convolve}, which does not detect separability, takes \(2.5\) ms and
\(53.2\) ms, while its dedicated \mcode{gaussian\_filter} takes \(1.6\) ms and
\(3.3\) ms (\texttt{docs/spec/builtins/image-processing.md}, ``Separable
kernels''). In Mathilda any rank-one kernel gets the fast path automatically.}

\subsection{Large kernels and the Fourier transform}

A kernel that is neither separable nor small is a candidate for the fast Fourier
transform, by the convolution theorem. A \emph{motion blur}\index{motion blur}
is a good example. Averaging along a diagonal line, \mcode{IdentityMatrix[61]/61.},
has rank \(61\), so it cannot be factored:

\plotcell[0.8\linewidth]{image-processing/motion}

\calloutnote{Under the hood}{The dispatcher in \texttt{src/imagefilter.c} tries
three methods in order: the separable factorisation, then the transform if a cost
model predicts it is cheaper, then the direct loop. A transform computes
\emph{circular} convolution, which would wrap the right edge of the image into
the left. To prevent that, the image is first padded by replication to
\(h + k_h - 1\) rows and \(w + k_w - 1\) columns, so no output pixel can read a
wrapped value, and the answer is read from a shifted window of the result.
Measured at \(512\times512\), the transform costs about \(2.7\) ms whatever the
kernel size, while the direct loop costs \(10.6\) ms for \(9\times9\) and
\(199\) ms for \(32\times32\). The crossover is set at \(7\times7\). Below it the
direct loop is kept, so small kernels stay bit-exact.}

\subsection{Sharpening}
\label{ssec:img-sharpen}

Sharpening is the reverse of smoothing: it amplifies the high frequencies that a
blur attenuated. A standard sharpening kernel is
\[
  \begin{pmatrix} 0 & -1 & 0 \\ -1 & 5 & -1 \\ 0 & -1 & 0 \end{pmatrix}
  = \delta - \nabla^2_{\text{discrete}},
\]
the identity minus the five-point discrete Laplacian. It leaves flat regions
unchanged (the entries sum to one) and exaggerates any difference between a pixel
and its neighbours. Applied to a blurred image, it restores some of the edge
contrast:

\plotcell[0.8\linewidth]{image-processing/sharpen}

\calloutnote{Pitfall}{Sharpening cannot restore information that the blur
destroyed. It steepens the transitions that remain, and it amplifies noise at the
same time, because noise is high-frequency too. Results above \(1\) and below
\(0\) are kept in the real image and are only clipped when the image is
exported.}

\section{Rank filters and noise}
\label{sec:img-rank}

A \emph{rank filter}\index{rank filter} sorts the values in each window and
takes one of them. \B{MedianFilter} takes the middle one. A median is not a
weighted sum, so it is not a convolution, and it behaves differently on outliers:

\begin{codepairs}
\pair{image-processing/rank/1}{One bright outlier in a flat grey field.}
\pair{image-processing/rank/2}{The median removes it completely.}
\pair{image-processing/rank/3}{The mean spreads it over the whole window.}
\pair{image-processing/rank/4}{A \(3\times3\) window of the numbers
\(1..9\).}
\pair{image-processing/rank/5}{Its median is \(5\). The median of the three row
medians is \(4\).}
\pair{image-processing/rank/6}{\mcode{MedianFilter} gives the true median.}
\end{codepairs}

Lines 5 and 6 show that the median is not separable. A sum, a maximum, and a
minimum over a rectangle can each be computed row by row and then column by
column, because they do not depend on how the values are grouped. A median depends
on each value's rank within the whole window, and grouping loses that
information. The shortcut in line 5 is fast and wrong.

The practical case is \emph{salt-and-pepper noise}\index{noise!salt-and-pepper},
where a fraction of pixels are forced to black or white. Below, from the top: the
noisy image, a \(3\times3\) mean filter, and a \(3\times3\) median filter. The
mean smears each outlier into a grey blotch. The median removes almost all of them
and keeps the edges sharp.

\plotcell[0.4\linewidth]{image-processing/median}

\calloutnote{Performance}{The window is sorted by insertion sort. For the radii
used in practice (\(1\) to \(3\), windows of \(9\) to \(49\) values) a small
contiguous sort beats a histogram or running median on constant factors. At
\(512\times512\), radius \(1\) takes \(3.13\) ms against SciPy's \(6.8\) ms, and
radius \(2\) takes \(10.27\) ms against \(15.3\) ms. The cost grows with the
window, so at large radii a histogram method would win.}

\section{Derivatives, edges, and corners}
\label{sec:img-edges}

\subsection{Derivative filters}

An edge is where brightness changes quickly, and edge detection therefore
starts with derivatives. The filter \B{DerivativeFilter} with orders \(\{n, m\}\) takes the \(n\)th
derivative down the rows and the \(m\)th across the columns. Each is a separable
kernel built from three one-dimensional stencils:
\[
  \text{order } 0:\ \tfrac14(1, 2, 1), \qquad
  \text{order } 1:\ \tfrac12(-1, 0, 1), \qquad
  \text{order } 2:\ (1, -2, 1).
\]
The smoothing stencil of order \(0\) is applied along the other axis, so order
\(\{0, 1\}\) is the Sobel operator for \(\partial/\partial x\).

\begin{codepairs}
\pair{image-processing/derivs/1}{A ramp rising by \(0.1\) per pixel.}
\pair{image-processing/derivs/2}{The \(x\) derivative is the slope, \(0.1\). The
ends are halved because the replicated border is flat.}
\pair{image-processing/derivs/3}{Nothing changes down the rows.}
\pair{image-processing/derivs/4}{A vertical step from \(0\) to \(1\).}
\pair{image-processing/derivs/5}{The central difference spreads the jump over two
pixels.}
\pair{image-processing/derivs/6}{\B{GradientFilter} is the magnitude
\(\sqrt{f_x^2 + f_y^2}\).}
\end{codepairs}

\calloutnote{Pitfall}{The textbook Sobel kernel has integer entries
\((1,2,1)\otimes(-1,0,1)\) and reports a gradient eight times the true slope.
Mathilda's stencils are normalised, so a ramp of slope \(0.1\) reports \(0.1\). The
ranking of edges is the same either way, but thresholds taken from a textbook
that uses the integer kernel must be divided by eight.}

The gradient magnitude is used, and not \(|f_x| + |f_y|\), because the magnitude
does not depend on the edge's orientation. A diagonal edge shows the difference:

\begin{codepairs}
\pair{image-processing/derivs/8}{A step along the diagonal.}
\pair{image-processing/derivs/9}{Both partial derivatives have size \(0.375\) at
the edge.}
\pair{image-processing/derivs/10}{The magnitude is \(0.53\). The absolute sum would
report \(0.75\), a factor \(\sqrt2\) larger, only because the edge is at
\(45^\circ\).}
\end{codepairs}

On a picture, the \(x\) derivative (top) responds to vertical edges with a sign:
white where brightness rises to the right, black where it falls. Zero is mid-grey,
after \mcode{ImageAdjust} has mapped the signed range onto \([0,1]\) for display.
The \(y\) derivative (middle) responds to horizontal edges. The magnitude (bottom)
is the outline in every direction.

\plotcell[0.4\linewidth]{image-processing/derivatives}

\subsection{Canny edge detection}

\B{EdgeDetect} implements the \emph{Canny detector}\index{Canny edge detector}. It
has four stages, and each one fixes a failure of the one before:

\begin{enumerate}
\item \textbf{Smooth.} Differencing amplifies noise, so the image is first blurred
with a Gaussian, radius \(2\) by default.
\item \textbf{Gradient.} Magnitude and direction come from the normalised Sobel
pair.
\item \textbf{Non-maximum suppression.} The magnitude forms a ridge several pixels
wide. A pixel is kept only if it is a maximum along the gradient direction, which
thins the ridge to one pixel.
\item \textbf{Hysteresis.} A single threshold either breaks weak stretches of a
long edge or admits noise. Canny uses two. A pixel above the high threshold is an
edge, and a pixel above the low threshold is kept if it is connected to one.
\end{enumerate}

\begin{codepairs}
\pair{image-processing/derivs/7}{With no smoothing, the step in the previous
session gives an edge exactly one pixel wide.}
\end{codepairs}

\calloutnote{Under the hood}{The one-pixel width depends on a detail. The central
difference of a clean step is an even plateau, \(0.5\) at two adjacent pixels (line
5 above). A test of the form ``at least both neighbours'' keeps both and draws a
two-pixel edge. \texttt{src/imagefilter.c} compares strictly against the backward
neighbour and loosely against the forward one, so ties resolve to one side. The
high threshold defaults to Otsu's threshold (Section~\ref{sec:img-threshold}) on
the \emph{suppressed} magnitude, and the low threshold is \(0.4\) times the high.
scikit-image's \mcode{canny} draws the same synthetic step two pixels wide and
discards border rows. At \(512\times512\) Mathilda takes \(9.1\) ms and
scikit-image \(19.9\) ms.}

The smoothing stage matters most on noisy data. Below, from the top: an image
with uniform noise added, Canny with no smoothing, and Canny with smoothing
radius \(4\). Without smoothing, the noise produces edges everywhere.

\plotcell[0.4\linewidth]{image-processing/canny}

\subsection{Corners}

A corner is a point where the brightness changes in two independent directions.
The standard tool is the \emph{structure tensor}\index{structure tensor}
\[
  M = \begin{pmatrix}
    \langle f_x^2 \rangle & \langle f_x f_y \rangle \\
    \langle f_x f_y \rangle & \langle f_y^2 \rangle
  \end{pmatrix},
\]
where \(\langle\cdot\rangle\) is a Gaussian-weighted average over a window. Both
eigenvalues of \(M\) are small on a flat region. On a straight edge every
gradient in the window is parallel, \(M\) has rank one, and one eigenvalue is
zero. At a corner both are large. \B{CornerFilter} gives the smaller eigenvalue
\(\lambda_{\min}\) (the Shi--Tomasi response) by default, or Harris's
\(\det M - 0.04 \operatorname{tr}(M)^2\). \B{ImageCorners} returns the positions of
the strongest local maxima:

\begin{codepairs}
\pair{image-processing/corners/1}{A bright square.}
\pair{image-processing/corners/2}{Its four corners, each given as a row and a
column.}
\pair{image-processing/corners/3}{A straight vertical edge.}
\pair{image-processing/corners/4}{Its Shi--Tomasi response is exactly zero
everywhere.}
\pair{image-processing/corners/5}{The Harris response goes negative on the
edge.}
\pair{image-processing/corners/6}{A checkerboard of \(6\times6\) squares.}
\pair{image-processing/corners/7}{Its nine interior corners.}
\pair{image-processing/corners/8}{Rotating the image rotates the response. The
two orders agree to rounding.}
\end{codepairs}

The response of the shapes image is bright at the four corners of the rectangle.
The circle also responds, in a dotted ring. A digitised circle is a staircase of
tiny steps, and each step is a small corner.

\plotcell[0.7\linewidth]{image-processing/corners}

\calloutnote{Pitfall}{\mcode{ImageCorners} returns \(\{\text{row},
\text{column}\}\) positions that index \mcode{ImageData} directly. Mathematica
returns \(\{x, y\}\) measured from the bottom left. The featured example on
template matching converts between the two.}

\calloutnote{Under the hood}{\(\lambda_{\min}\) is computed as
\(\tfrac12\bigl(S_{xx} + S_{yy} - \sqrt{(S_{xx}-S_{yy})^2 + 4S_{xy}^2}\bigr)\). The
form \(\sqrt{\operatorname{tr}^2 - 4\det}\) is algebraically equal, but
cancellation can make its argument slightly negative and produce a NaN.
\mcode{ImageCorners} applies \(3\times3\) non-maximum suppression, then the
threshold (a fraction of the largest response, \(0.05\) by default), then an
optional minimum separation, choosing greedily from the strongest down. At
\(512\times512\) the positions take \(6.01\) ms against \(29.8\) ms for
scikit-image's \mcode{corner\_shi\_tomasi} with \mcode{corner\_peaks}.}

\subsection{Template matching}
\label{ssec:img-ncc}

To find a small pattern \(t\) inside an image, slide it across and score each
position. Plain correlation is a poor score because it rewards brightness: a white
patch beats a correct but darker match. The \emph{normalised cross-correlation}
\index{normalised cross-correlation} subtracts the mean of each window and divides
by the standard deviations,
\[
  \operatorname{NCC}(y, x) =
  \frac{\sum (f - \bar f_{w})(t - \bar t)}
       {\sqrt{\sum (f - \bar f_{w})^2 \sum (t - \bar t)^2}},
\]
where the sums run over the window \(w\) at \((y, x)\). The score lies in
\([-1, 1]\), and it is unchanged if the image is multiplied by a positive constant
or has a constant added:

\begin{codepairs}
\pair{image-processing/ncc/1}{A noise image, \(40\) wide and \(30\) high.}
\pair{image-processing/ncc/2}{The template is a \(7\times7\) crop centred at row
\(14\), column \(24\).}
\pair{image-processing/ncc/3}{Score every position.}
\pair{image-processing/ncc/4}{The best match is at the crop's centre, with a
score of \(1\).}
\pair{image-processing/ncc/5}{Paint a white block into the lower left.}
\pair{image-processing/ncc/6}{Then change the contrast and offset of the
whole image.}
\pair{image-processing/ncc/7}{Plain correlation is computed.}
\pair{image-processing/ncc/8}{It now peaks inside the white block.}
\pair{image-processing/ncc/9}{The normalised score is computed.}
\pair{image-processing/ncc/10}{It still finds the true match.}
\end{codepairs}

\calloutnote{Under the hood}{The definition needs the window sum and the window
sum of squares at every position. Using the identities
\(\sum (f - \bar f)(t - \bar t) = \sum f t - \bar t \sum f\) and
\(\sum (f - \bar f)^2 = \sum f^2 - (\sum f)^2/m\), with \(m\) the number of pixels
in the window, Mathilda reads both from \emph{summed-area tables}\index{summed-area
table}. Each window sum then costs four lookups whatever the template size. The
remaining cross term \(\sum f t\) is an ordinary correlation and goes through the
Fourier path. At \(512\times512\) with a \(32\times32\) template, matching takes
\(3.55\) ms against \(146.5\) ms for the same algorithm in NumPy and SciPy.}

\calloutnote{Pitfall}{The template is passed as a matrix of samples, as in line 2,
not as an \mcode{Image}. A window with no variation has no shape to compare and
scores \(0\).}

\section{Thresholding}
\label{sec:img-threshold}

\emph{Binarisation}\index{binarisation} turns a grey image into a bit image,
separating foreground from background. \B{Binarize} with a number \(t\) sets
every pixel strictly above \(t\) to \(1\). With no number, it chooses \(t\) by
\emph{Otsu's method}\index{Otsu's method}, which \B{FindThreshold} exposes on its
own.

\begin{theorem}[Otsu, 1979]
Split the pixels at a threshold \(t\) into two classes with weights \(w_0, w_1\)
(the fractions of pixels in each) and means \(\mu_0, \mu_1\). Minimising the
weighted within-class variance \(w_0\sigma_0^2 + w_1\sigma_1^2\) is equivalent to
maximising the between-class variance
\[
  \sigma_B^2(t) = w_0\, w_1\, (\mu_0 - \mu_1)^2 .
\]
\end{theorem}

The equivalence holds because the total variance is the sum of the two, and the
total does not depend on \(t\). The between-class form needs only running sums,
so one pass over a \(256\)-bin histogram finds the best threshold.

\begin{codepairs}
\pair{image-processing/threshold/1}{Two clusters of values, around \(0.15\) and
\(0.85\).}
\pair{image-processing/threshold/2}{Otsu's threshold falls between the clusters,
at a histogram bin level.}
\pair{image-processing/threshold/3}{The two clusters separate.}
\pair{image-processing/threshold/4}{A pixel exactly at the threshold becomes
\(0\): the comparison is strict.}
\pair{image-processing/threshold/5}{The result is a bit image.}
\end{codepairs}

A single threshold fails when the lighting varies across the image. If one side
of a page is darker than the other, no one number separates ink from paper on both
sides. \B{LocalAdaptiveBinarize}\mcode{[img, r, \{c1, c2, c3\}]} compares each pixel
with its own neighbourhood instead. The threshold at each pixel is
\[
  T(y, x) = c_1\, \mu_r(y, x) + c_2\, \sigma_r(y, x) + c_3 ,
\]
where \(\mu_r\) and \(\sigma_r\) are the mean and standard deviation over the
\((2r+1)\times(2r+1)\) window. The default \(\{1, 0, 0\}\) is Bradley's method,
and a negative \(c_2\) gives Sauvola's.

\begin{codepairs}
\pair{image-processing/threshold/6}{A checkerboard under a strong lighting ramp.}
\pair{image-processing/threshold/7}{Otsu's global threshold turns the left side
entirely black.}
\pair{image-processing/threshold/8}{The local threshold keeps both tones
there.}
\end{codepairs}

Below, from the top: the unevenly lit checkerboard, Otsu's global threshold, and
the local threshold at radius \(32\).

\plotcell[0.7\linewidth]{image-processing/binarize}

\calloutnote{Performance}{The window mean and variance come from summed-area
tables, so the cost per pixel does not depend on \(r\). At \(512\times512\),
radius \(2\) takes \(0.58\) ms and radius \(32\) takes \(0.73\) ms, against about
\(1.7\) ms for scikit-image's \mcode{threshold\_local} at every radius. An earlier
version built the result as \(262{,}144\) separate integer expressions, and that
construction took longer than the thresholding. Emitting a packed buffer instead
also made global \mcode{Binarize} \(23\times\) faster, from \(6.47\) ms to
\(0.278\) ms.}

\calloutnote{Pitfall}{With the default \(\{1, 0, 0\}\) the threshold \emph{is} the
local mean, and on a smooth ramp a pixel can equal its local mean. The decision on
such a pixel then depends on the last bit of the arithmetic. A small offset, such
as \(c_3 = -0.05\), moves the threshold away from the tie. The featured example on
page binarisation uses one.}

\section{Mathematical morphology}
\label{sec:img-morphology}

\emph{Mathematical morphology}\index{mathematical morphology} studies shapes by
probing them with a small set, the \emph{structuring element} \(B\). For a grey
image \(f\), the two basic operations are
\[
  (f \oplus B)(p) = \max_{b \in B} f(p + b), \qquad
  (f \ominus B)(p) = \min_{b \in B} f(p + b).
\]
\B{Dilation} (\(\oplus\)) takes the maximum over the element, so bright regions
grow. \B{Erosion} (\(\ominus\)) takes the minimum, so they shrink. On a bit
image these are the Minkowski sum and difference of the foreground with \(B\).
\B{Opening} erodes and then dilates. It removes bright features smaller than
\(B\) and restores the rest. \B{Closing} dilates and then erodes. It fills dark
features smaller than \(B\).

\begin{codepairs}
\pair{image-processing/morph/1}{A single bright pixel.}
\pair{image-processing/morph/2}{Dilated by a cross, it becomes the cross. The
output of a dilated point \emph{is} the element.}
\pair{image-processing/morph/3}{Only the element's support matters. A radius
\(1\), a \mcode{BoxMatrix}, and a box of fives give the same result.}
\pair{image-processing/morph/4}{A random grey image for testing the laws.}
\pair{image-processing/morph/5}{Duality: erosion is dilation of the negative,
negated. The error is exactly zero.}
\pair{image-processing/morph/6}{Four operations are run at radius \(2\).}
\pair{image-processing/morph/7}{They bracket the image pixel by pixel: erosion
\(\le\) opening \(\le f \le\) closing \(\le\) dilation.}
\pair{image-processing/morph/8}{Opening is idempotent. Opening twice is opening
once.}
\end{codepairs}

These laws are the right tests of an implementation because each fails for a
different bug. Duality fails if the maximum and minimum are swapped or if the
border rule treats bright and dark differently. The ordering fails if the element
is off-centre. Idempotence fails if the pair is composed wrongly. Replicate padding
is what makes all three hold at the border as well as in the interior.

A speckled disk shows all four operations. Top row: the original, erosion, and
dilation. Bottom row: the original, opening, and closing. Erosion eats into the
disk through every hole. Dilation grows every speck of noise. Opening removes the
isolated specks outside the disk, and closing fills the holes inside it.

\plotcell[0.8\linewidth]{image-processing/morphology}

\calloutnote{Theory}{The maximum over a window of width \(k\) can be computed with
three comparisons per pixel whatever \(k\) is. This is the \emph{van
Herk--Gil--Werman} algorithm\index{van Herk--Gil--Werman algorithm}. Cut the line
into blocks of length \(k\). Within each block compute a running maximum forwards
and another backwards. Any window of length \(k\) overlaps exactly two adjacent
blocks, so its maximum is the larger of the backward maximum at its start and the
forward maximum at its end. Because the maximum is associative, the result is
exact. A square element is separable, so the same idea handles two dimensions one
axis at a time.}

\calloutnote{Performance}{With van Herk--Gil--Werman, dilation of a
\(512\times512\) image costs about \(2.0\) ms at every radius from \(2\) to \(32\),
against \(2.3\) to \(2.4\) ms for SciPy's \mcode{grey\_dilation}
(\texttt{docs/spec/builtins/image-processing.md}, ``Morphology'').}

\calloutnote{Pitfall}{In Mathilda the morphological operations return a
\mcode{"Real"} image even when the input is a bit image. The values are still
exactly \(0\) and \(1\), but \mcode{ImageData[result, "Bit"]} is refused because
\mcode{"Bit"} is not the result's type. Read the pixels with plain
\mcode{ImageData}, or apply \mcode{Binarize} if a bit image is needed.}

\subsection{Connected components}

\B{MorphologicalComponents} labels the connected regions of the foreground. It
returns an integer matrix: \(0\) for the background, and \(1, 2, \dots\) for the
components in the order they are first met in a raster scan. Whether two pixels
touching only at a corner are connected is a choice, and the option
\mcode{CornerNeighbors} makes it:

\begin{codepairs}
\pair{image-processing/components/1}{With the default 8-connectivity, diagonal
neighbours join: one component.}
\pair{image-processing/components/2}{With 4-connectivity they are two.}
\pair{image-processing/components/3}{A U shape.}
\pair{image-processing/components/4}{Its two arms are one component, joined at
the bottom.}
\pair{image-processing/components/5}{Four isolated corner pixels.}
\pair{image-processing/components/6}{Four components, and one after dilation.
Dilation can merge components but never split them.}
\pair{image-processing/components/7}{A second argument sets the foreground
threshold.}
\end{codepairs}

\calloutnote{Under the hood}{The U shape is the case that separates a correct
labelling from a naive one. The scan moves left to right and top to bottom, and
at each pixel it can see only neighbours already visited. It meets the two arms
first and gives them different labels, and only the bottom row shows that they
are connected. Mathilda records such equivalences in a union-find structure during
the first pass and resolves them in a second pass, which also renumbers the labels
to \(1..k\) with no gaps.}

The result is a matrix of labels, which can be turned back into an image for
display. Here each label becomes its own grey level:

\plotcell[0.7\linewidth]{image-processing/components}

\calloutnote{Pitfall}{Labels are indices. Wrapping a label
matrix in \mcode{Image} unchanged, when its values run up to \(255\), gives a
byte image, and \mcode{ImageData} would then divide every label by \(255\). The
figure's input rescales the labels explicitly first. When the function landed, labelling a \(512\times512\)
image took \(5.2\) ms against SciPy's \(1.2\) ms. Most of that time was spent
building the output matrix as expressions. The labelling itself was fast.}

\subsection{The distance transform}

\B{DistanceTransform} replaces each foreground pixel by its distance to the
nearest background pixel. Mathilda computes the \emph{exact} Euclidean distance:

\begin{codepairs}
\pair{image-processing/distance/1}{A field of foreground with one background pixel
in the top-left corner.}
\pair{image-processing/distance/2}{Compute the transform.}
\pair{image-processing/distance/3}{Three columns across and four rows down is at
distance exactly \(5\).}
\pair{image-processing/distance/4}{The diagonal neighbour is at \(\sqrt 2\).}
\pair{image-processing/distance/5}{The top-left corner of the result.}
\pair{image-processing/distance/6}{A small disk.}
\pair{image-processing/distance/7}{A pixel survives one \(3\times3\) erosion
exactly when its distance is at least \(2\).}
\end{codepairs}

\calloutnote{Theory}{The classic two-pass \emph{chamfer} transform propagates
integer step costs and cannot represent \(\sqrt2\), so its diagonal distances are
a few percent off. Mathilda uses the method of Felzenszwalb and Huttenlocher. The
squared Euclidean distance is a sum over the axes, so it can be minimised one axis
at a time. Along one axis the problem is
\(D(x) = \min_y \bigl((x - y)^2 + g(y)\bigr)\), the lower envelope of a family of
parabolas that all have the same shape. Any two such parabolas cross exactly once,
so the envelope can be built in a single sweep with a stack, in linear time. The
square root is taken once, at the end.}

Rendered after \mcode{ImageAdjust}, the transform of the two shapes brightens
toward their centres. The ridge lines in the rectangle are its \emph{medial axis},
the points equidistant from two sides.

\plotcell[0.7\linewidth]{image-processing/distance}

\calloutnote{Performance}{On the same \(512\times512\) mask, the transform takes
\(2.86\) ms against \(7.3\) ms for SciPy's \mcode{distance\_transform\_edt}, and
the two agree exactly: both report a maximum of \(6\) and a total of
\(323{,}895\).}

\subsection{Thinning and pruning}

\B{Thinning} reduces a shape to its \emph{skeleton}\index{skeleton}, a curve one
pixel wide that keeps the shape's connectivity. \B{Pruning} then shortens the free
ends of the skeleton, which removes the short spurs that thinning grows at small
irregularities of the boundary.

\begin{codepairs}
\pair{image-processing/thin/1}{A bar three pixels thick.}
\pair{image-processing/thin/2}{Its skeleton is the centre line.}
\pair{image-processing/thin/3}{Pruning twice removes two pixels from each end.}
\pair{image-processing/thin/4}{An isolated pixel has no free end, so pruning
leaves it alone.}
\pair{image-processing/thin/5}{A grey image is thresholded at \(0.5\) first.}
\end{codepairs}

\calloutnote{Under the hood}{\texttt{src/imagethin.c} uses the
\emph{Zhang--Suen}\index{Zhang--Suen thinning} algorithm. A pixel may be deleted
if it has between two and six foreground neighbours and exactly one transition from
background to foreground as you walk once around its eight neighbours. The two
subiterations of each pass add different conditions on the north, south, east,
and west neighbours, so they remove pixels from opposite sides in turn. Deleting
every removable pixel in a single pass would cut diagonal lines, because two
diagonal neighbours can each be removable while removing both disconnects the
shape. The passes repeat until one deletes nothing.}

A thick ring with a spur thins to a one-pixel circle with a whisker, and ten rounds
of pruning remove the whisker:

\plotcell[0.8\linewidth]{image-processing/thinning}

\section{Composition}
\label{sec:img-compose}

\B{ImageCompose} lays one image over another. It alpha-composites the overlay onto
the base, keeps the base's size, and clips whatever falls outside.
\B{ImageAssemble} tiles a grid of images into one, which is how the figures in
this chapter place results side by side.

\begin{codepairs}
\pair{image-processing/compose/1}{A black \(6\times4\) base.}
\pair{image-processing/compose/2}{A white \(2\times2\) patch.}
\pair{image-processing/compose/3}{By default the overlay is centred.}
\pair{image-processing/compose/4}{A position \(\{x, y\}\) places the overlay's
centre, with \(y\) counted from the \emph{bottom}. An opacity of \(0.5\) gives
grey.}
\pair{image-processing/compose/5}{Composing a colour overlay onto a grey base
gives a colour result.}
\pair{image-processing/compose/6}{Tiles keep their natural sizes. A gap is left
blank.}
\end{codepairs}

Compositing uses the standard \emph{over} operator\index{alpha compositing}: with
overlay colour \(c_o\) and opacity \(\alpha\) on a base colour \(c_b\), the result
is \(\alpha c_o + (1 - \alpha) c_b\). Three half-transparent disks placed on a
gradient show the background through them:

\plotcell[0.6\linewidth]{image-processing/compose}

\calloutnote{Pitfall}{\mcode{ImageCompose} takes positions in Mathematica's image
coordinates, \(x\) from the left and \(y\) from the bottom, while \mcode{ImageData}
and \mcode{ImageCorners} count rows from the top. Converting a row \(r\) of an
image of height \(h\) gives \(y = h - r + 1\). A composition that got this wrong
would look right on any test image that is symmetric top to bottom.}

\calloutnote{Under the hood}{\texttt{src/imagecompose.c} decides how to combine
channel counts in one place. A grey image next to a colour image is replicated into
three equal channels, never padded with zeros, because grey means the same value in
every channel and zero-padding would turn a grey pixel red.}

\section{Volumes}
\label{sec:img-volumes}

\B{Image3D} is the three-dimensional version of \mcode{Image}. Its data is a stack
of slices, depth by height by width, so it is indexed \mcode{[[z, y, x]]}. CT and
MRI scans, confocal microscopy, and simulation output all come in this form.

\begin{codepairs}
\pair{image-processing/volume/1}{Two slices of three rows of four voxels. The
value encodes the position.}
\pair{image-processing/volume/2}{\mcode{ImageDimensions} is fully reversed. It
lists width, then height, then depth.}
\pair{image-processing/volume/3}{A volume is not a plane. \mcode{ImageQ} is false
and \B{Image3DQ} is true.}
\pair{image-processing/volume/4}{Slice \(2\), row \(3\), column \(4\) holds
\(0.234\).}
\pair{image-processing/volume/5}{A solid ball of radius \(5\).}
\pair{image-processing/volume/6}{Its voxel sum, and after a Gaussian, an erosion,
and a dilation. The Gaussian preserves the total exactly.}
\pair{image-processing/volume/7}{The distance transform works in three
dimensions: the centre is \(5\) voxels from the outside.}
\pair{image-processing/volume/8}{Padding adds two voxels on every side of every
axis.}
\pair{image-processing/volume/9}{\mcode{Front} reflects along the depth axis, so
the first slice is now the old second slice.}
\end{codepairs}

With three axes there are six possible orderings, and a cubic test volume cannot
tell any of them apart. That is why the session uses extents \(4, 3, 2\).

A slice of a volume is an ordinary image, so \mcode{Part} and \mcode{Image} are
enough to look inside one. Below are three slices through a ball of radius \(10\)
(top row), the same slices after a radius-\(2\) dilation (middle), and after a
radius-\(3\) Gaussian blur (bottom). Each operation acts in all three directions,
so the slice nearest the edge of the ball (left column) changes the most.

\plotcell[0.45\linewidth]{image-processing/volume}

\calloutnote{Performance}{Separability pays more in three dimensions. A separable
kernel costs \(3(2r+1)\) multiplications per voxel instead of \((2r+1)^3\): at
radius \(4\), \(27\) instead of \(729\). On a \(64^3\) volume a radius-\(4\)
Gaussian takes \(2.92\) ms. SciPy's general \mcode{ndimage.convolve} takes
\(155.1\) ms and its separable \mcode{gaussian\_filter} \(2.9\) ms
(\texttt{docs/spec/builtins/image-processing.md}, ``Volumetric convolution'';
\texttt{benchmarks/69-image3d-fft-convolve} through
\texttt{benchmarks/78-image3d-reflect} for the other volumetric heads).}

\section{Storage and speed}
\label{sec:img-speed}

An image operation does very little arithmetic per pixel. A \(3\times3\) blur is
nine multiply-adds. The cost that dominates is often the movement of data, and in
a system built on expressions the obvious representation, a nested list of
numbers, is expensive: each pixel is a separate heap object. Mathilda therefore
stores the pixels of an image in a machine buffer:

\begin{codepairs}
\pair{image-processing/storage/1}{A filtered image.}
\pair{image-processing/storage/2}{Its pixels are an \mcode{NDArray}, a flat
buffer of machine doubles.}
\pair{image-processing/storage/3}{\mcode{ImageData} hands them out as an ordinary
list.}
\pair{image-processing/storage/4}{Pixels given to \mcode{Image} as a list of
machine numbers are stored as a buffer too.}
\end{codepairs}

\calloutnote{Under the hood}{Mathilda has two array surfaces (see
Section~\ref{sec:ds-packed}): a \emph{packed list}, which looks like an ordinary
\mcode{List}, and a visible \mcode{NDArray}. Images use the visible form. The
evaluator unpacks a packed list that is still at rest when evaluation finishes,
which is correct for an ordinary function call but would unpack an image's pixels
on every evaluation. A visible \mcode{NDArray} is never unpacked that way. The
image heads are also listed as packed-aware in \texttt{src/pack.c}, so a buffer
passed between filters is never converted into expressions on the way in.}

The effect of the change is large, and it is measured
(\texttt{docs/spec/builtins/image-processing.md}, ``Storage: machine buffers'').
On a \(512\times512\) image, the fixed cost of getting pixels into and out of a
filter fell from \(4.62\) ms to \(0.98\) ms, and a chain of three
\mcode{GaussianFilter} calls fell from \(17.3\) ms to \(4.43\) ms.
Table~\ref{tab:img-bench} collects the comparisons quoted in this chapter. All
times are for a \(512\times512\) grey image unless stated, on Apple silicon.

\begin{table}[ht]
\centering
\small
\begin{tabular}{@{}lrlr@{}}
\toprule
Operation & Mathilda & Reference & Reference time \\
\midrule
Gaussian, radius 8 (separable) & 3.44 & SciPy \texttt{gaussian\_filter} & 3.3 \\
Gaussian, radius 8 (separable) & 3.44 & SciPy \texttt{convolve} & 53.2 \\
Resize \(512\to1024\), bilinear & 2.76 & scikit-image \texttt{resize} & 14.3 \\
Median, radius 1 & 3.13 & SciPy \texttt{ndimage} median & 6.8 \\
Dilation, radius 16 & 2.03 & SciPy \texttt{grey\_dilation} & 2.3 \\
Distance transform & 2.86 & SciPy \texttt{distance\_transform\_edt} & 7.3 \\
Canny edges & 9.1 & scikit-image \texttt{canny} & 19.9 \\
Corner positions & 6.01 & scikit-image \texttt{corner\_peaks} & 29.8 \\
Template match, \(32\times32\) & 3.55 & NumPy/SciPy, same algorithm & 146.5 \\
Local threshold, radius 16 & 0.66 & scikit-image \texttt{threshold\_local} & 1.68 \\
Contrast stretch & 0.76 & scikit-image \texttt{rescale\_intensity} & 7.7 \\
Component labelling & 5.20 & SciPy \texttt{label} & 1.2 \\
Gaussian on \(64^3\), radius 4 & 2.92 & SciPy \texttt{gaussian\_filter} & 2.9 \\
\bottomrule
\end{tabular}
\caption{Image operations in milliseconds. Sources:
\texttt{docs/spec/builtins/image-processing.md} and the benchmark pairs under
\texttt{benchmarks/} that it cites. The labelling figure is from when that function
landed and predates later packing work.}
\label{tab:img-bench}
\end{table}

Two rows need comment. The Gaussian is at parity with SciPy's dedicated separable
routine, and \(15\times\) faster than its general convolution only because that
routine does not detect separability. Component labelling is the one row where
Mathilda was slower. The labelling itself took about \(0.6\) ms. The rest was the
construction of the \(262{,}144\)-entry label matrix as expressions, the same
bottleneck that the buffer storage removed for images.

\calloutnote{Performance}{Three times in this subsystem an operation ran \(4\) to
\(23\) times slower than it should have while giving correct answers, because the
pixels were being converted between buffers and expressions. No correctness test
can catch that. \mcode{make check-image-packing} now checks mechanically that every
image-returning head hands back a packed buffer, at both two and three
dimensions.}

The image heads are containers and are not part of the subset that
\B{Compile} lowers to machine code:

\begin{codepairs}
\pair{image-processing/storage/5}{A compiled body that calls an image filter falls
back to the interpreter.}
\end{codepairs}

This costs little, because each image head is already a single call into compiled
C that processes the whole buffer. Put the image operations outside a compiled
function, and compile only the per-pixel arithmetic of your own that surrounds
them.

\section{Limitations}
\label{sec:img-limits}

The image subsystem covers the core of a first course in image processing. It
differs from Mathematica's in several ways that you will meet quickly:

\begin{itemize}
\item \mcode{ColorConvert} converts only to grey. There are no HSB, Lab, or other
colour spaces yet.
\item Images do not support arithmetic directly. \mcode{Image[\{\{0.2\}\}] +
Image[\{\{0.1\}\}]} stays unevaluated. Do the arithmetic on \mcode{ImageData}
and wrap the result in \mcode{Image}, as several figures in this chapter do.
\item There is no \mcode{ExampleData} test-image library and no
\mcode{Rasterize}, so test images come from formulas or from files.
\item \mcode{ImageRotate} does not accept volumes, since a rotation in three
dimensions needs an axis.
\item The template of \mcode{ImageCorrelate} is a matrix, not an
\mcode{Image}.
\item \mcode{ImageCorners} returns \(\{\text{row}, \text{column}\}\), not
Mathematica's \(\{x, y\}\).
\item Morphological operations return real images even for bit input.
\item Some Mathematica image functions (segmentation by watershed, feature
descriptors, image alignment) have no counterpart yet.
\end{itemize}

\section{Featured examples}
\label{sec:img-featured}

The remaining pages are short worked tasks. Each takes one to three lines of real
work, and each builds its own synthetic input so that you can reproduce it
exactly. The input is a stand-in for the real thing: a photograph of coins, a
scanned page, a radiograph. The operations are the ones you would apply to the
real image.

\subsection*{Count the coins on a tray}

Twelve coins lie on a ribbed tray. Otsu's threshold separates metal from tray,
and counting the connected components counts the coins. The whole pipeline is the
third line:

\plotcell[0.55\linewidth]{image-processing/feat-coins}

\subsection*{Separate touching cells}

Cells in a micrograph often touch, and touching cells merge into one component.
Here nine discs merge into six. The distance transform is large only near the
centre of each disc, so keeping the pixels more than \(30\) from the background
leaves one small core per cell, and the count is right again. From the left: the
cells, their distance transform, and the cores.

\plotcell[\linewidth]{image-processing/feat-cells}

This works while the overlap between two cells is smaller than their radius.
Cells that overlap more than that need a watershed segmentation, which Mathilda
does not yet provide.

\subsection*{Binarise a page under uneven light}

A page photographed under a desk lamp is bright on one side and dark on the other.
Otsu's global threshold (middle) loses the whole dark half. A local threshold with
a small offset (bottom) recovers every glyph. On this image the recovered ink
matches the true ink mask pixel for pixel.

\plotcell[0.7\linewidth]{image-processing/feat-page}

\subsection*{Skeletonise a pen stroke}

Handwriting recognition often works on the centre line of a stroke. \mcode{Thinning} gives the centre line. \mcode{Pruning} then removes
spurs at the ends. The dot of the ``i'' thins to a single pixel and survives the
pruning, because an isolated pixel has no free end.

\plotcell[0.7\linewidth]{image-processing/feat-stroke}

\subsection*{Outline a bone in a noisy radiograph}

A long bone in cross-section is a bright shell of cortex around darker marrow. The
radiograph is noisy, so the smoothing radius of the Canny detector is raised to
\(6\). Both the outer and the inner boundary of the cortex come out as clean
closed curves.

\plotcell[0.7\linewidth]{image-processing/feat-xray}

\subsection*{Find a patch in a texture and mark it}

Normalised cross-correlation finds where an \(11\times11\) patch came from in a
random texture. The peak is reported as \(\{\text{row}, \text{column}\}\).
\mcode{ImageCompose} counts \(y\) from the bottom, so the frame is placed at
\(\{\text{col}, 96 - \text{row} + 1\}\). The frame image is transparent except
for its border, so the texture shows through.

\plotcell[0.6\linewidth]{image-processing/feat-target}

\subsection*{Replace the background}

A mask of the foreground becomes the alpha channel of the image. Composited onto a
new background, only the masked region of the original remains. The hole in the
mask lets the new sky show through.

\plotcell[0.7\linewidth]{image-processing/feat-background}

\subsection*{Locate the corners of a calibration board}

Camera calibration starts by finding the inner corners of a chessboard. With a
minimum separation of \(32\) pixels, \mcode{ImageCorners} returns one position per
corner, all \(35\) of them. The marks are dilated so that they are visible at
print size.

\plotcell[0.8\linewidth]{image-processing/feat-chess}

\subsection*{Sharpen with an unsharp mask}

The \emph{unsharp mask}\index{unsharp mask} is the photographer's sharpening
method: subtract a blurred copy from the image and add the difference back,
\(f + (f - G * f) = (2\delta - G) * f\). Written as one kernel, it is a single
convolution. The rings of the softened zone plate regain their contrast.

\plotcell[0.8\linewidth]{image-processing/feat-unsharp}

\subsection*{Make a six-colour poster}

Screen printing uses a few flat inks. Median cut chooses six colours that cover
the spread of a smooth sunset-coloured field, and every pixel takes the nearest of
them.

\plotcell[0.8\linewidth]{image-processing/feat-poster}

\subsection*{Look inside a scanned volume}

A CT scan is a volume. This one is a hollow shell with a small ball hidden
inside. A single slice (left) shows the ring of the shell and the ball as a disk.
The mean along the depth axis (right) simulates a projection radiograph, in which
the shell becomes a bright rim and the ball a faint shadow.

\plotcell[0.5\linewidth]{image-processing/feat-ctscan}

\subsection*{Where this connects}

Images draw on much of the rest of the book. Their pixels are the machine buffers
of Section~\ref{sec:ds-packed}, and \mcode{ImageData} converts them to the ordinary
lists that every list function accepts. The Fourier transform that speeds up large
convolutions is the one the numerical chapters use, and the convolution theorem is
the same theorem there. Random test images use the seeded generator of the rest of
the system, so every figure here can be reproduced. For drawing plots rather than
processing rasters, see Chapter~\ref{ch:graphics}, and for other file formats see
Chapter~\ref{ch:dataio}. For the options of any function in this chapter, type
\texttt{?Name} at the prompt to read its reference page.
